\section{Lean Mechanization and Verification Status}
\label{sec:mechanization}

The artifact uses Lean 4.31.0 without Mathlib. The current branch is
\texttt{research/modal-probabilistic-fine-graining}; the root module
\texttt{PEL4.lean} imports the relevant development. Appendix~\ref{app:formal}
provides the claim-to-declaration map, including the reliability refinement.
The explicit audit files are the machine-readable list of selected claims;
this selection is not a certification of every declaration in the repository.

\subsection{Verification contract}
A successful build establishes that declarations typecheck, not that all their
assumptions are proved. Version 0.4 therefore audits 52 manuscript declarations
in \texttt{PEL4/PaperAxiomAudit.lean} and 24 MPFG declarations in its separate
audit. The full build and both strict audits passed for source revision
\texttt{07fbc5193e0f7e43de565158c0eaf0cbb9c08dfb}; the review report records the
CI run and exact dependencies. Only \texttt{propext}, \texttt{Classical.choice},
and \texttt{Quot.sound}, or subsets thereof, occur in these selected chains.
\texttt{sorryAx}, project-specific axioms, and native-decision axioms are rejected.

The review found that some general headlines inherited native-evaluation
assumptions through elementary helpers. In eight source modules, the relevant
finite computations were changed from \texttt{native\_decide} to
\texttt{decide +kernel}, without changing statements or semantic definitions.
Legacy declarations outside the audited chains may still use native evaluation
or project axioms; no repository-wide axiom-freedom claim is made.

\subsection{Status discipline}
The manuscript follows the repository verification policy:
\begin{enumerate}[label=(\arabic*)]
\item \textbf{PROVED}: the stated theorem typechecks and its audited dependency
chain contains no project-specific unproved assumption. Explicit hypotheses
and model proof fields remain assumptions of the theorem's quantified statement.
\item \textbf{FINITE-WITNESS}: a finite construction or counterexample is
verified. Finite domain coverage is not a theorem about arbitrary model sizes.
\item \textbf{EXPERIMENTAL / INTERPRETIVE}: terminology, philosophical readings,
and prospective extensions without a corresponding established theorem.
\item \textbf{AXIOMATIC-PROTOTYPE}: a substantive property is postulated.
No manuscript headline is promoted from this category merely because it compiles.
\end{enumerate}
The proof sketches explain the main arguments for readers; Lean verifies their
formal counterparts. The geometry uses rational algebra, not topological
continuity. Numerical figure coordinates illustrate those statements and are
not additional model-realization certificates.

\paragraph{Reproducibility.}
From a clean clone, or after preserving local work:
\begin{verbatim}
git fetch origin
git switch research/modal-probabilistic-fine-graining
git pull --ff-only
lake build
python3 scripts/check_mpfg_axioms.py
python3 scripts/check_paper_axioms.py --report paper-axioms.json
python3 -m unittest discover -s scripts -p 'test_*.py'
cd paper
latexmk -pdf -interaction=nonstopmode -halt-on-error main.tex
\end{verbatim}
The mutable branch is for ongoing development; use a recorded commit for an
exact reproduction. Lean checking and manuscript rendering remain separate
checks. The committed PDF is the rendered working artifact, not evidence of
formal correctness by itself.
